Describe the organization or flow of the contents in this chapter before the start of the first section.
What can your readers learn from this chapter?
How did you organize your discussions and subsections?
You can add subsections but NOT sections.

Description of Subjects can be removed if this is not applicable to your research.

\section{Proposed Method}
Aside from the discussion regarding the proposed method, this must contain the block diagram(s) and possible flowchart(s).
Note that supporting circuits, block diagrams and flowcharts should be in the appendix.
For components, the order should be INPUT, PROCESS, OUTPUT.
The sequence should be as follows: system block diagram, flowchart(s), components, circuit(s), and final project.

DO NOT FORGET: 
All figures and tables should be mentioned in the paragraph or discussion.
Discuss the system block diagram in detail.
Each block should contain the process NOT the result of the process.
Variable names and the flow shown in the flowcharts should also be explained.
For components, the following should be included: picture of actual component; model number; technical specifications (only those that are significant for the project); purpose or function in the prototype; and reason for selection.
For the proposal hearing, CAD model with dimensions of the prototype should be shown and discussed.
After the final defense, include the actual picture at different angles and explain its step-by-step operation.

\section{Description of Experiments}
\label{sec: expts}
During the proposal period, describe in detail the possible experiments that will be performed.
For each experiment, it must have the objective(s), materials, and procedure.
Each subsection here should be an experiment.
All the details, settings, and assumptions should be discussed.
Tables may also be presented here.

\subsection{Title of Experiment of 1}
This experiment aims to measure something.
Note that each experiment should have a measureable objective that will help you determine/test a specific aspect/performance of your prototype.
You should VARY ONLY ONE VARIABLE PER EXPERIMENT.
The materials required are:
\begin{enumerate}
	\item material 1;
	\item material 2; and
	\item specify program name and include this in the appendix.
\end{enumerate}
The procedure is as follows:
\begin{enumerate}
	\item Each item here should be something to be done. 
	\item Descriptions or expected changes can be included RIGHT AFTER a specific procedure.
	\item For those that require tables, show the template of the table and specify in the procedure the column where the data should be inputted. 
	\item If a graph will be required, spcify the type of graph and data to be used.
	\item Lastly, instructions on processing the collected data should also be mentioned here.
\end{enumerate}
\subsection{Title of Experiment of 2}
This experiment aims to measure something.
The materials required are:
\begin{enumerate}
	\item material 1
	\item material 2
\end{enumerate}
The procedure is as follows:
\begin{enumerate}
	\item Load the program to something.
	\item Place one part in a corner.
\end{enumerate}

\section{Description of Subjects}
This section is optional or may be changed depending on your project.
For those who will use datasets, you can discuss the details here.